% !PRINTED 432
\weberpage{432}
%<<CONT
\[
\begin{array}{c|cccccccccccc}
I & 0 & 1 & 2 & 3 & 4 & 5 & 6 & 7 & 8 & 9 & 10 & 11 \\ \hline
N & 1 & 2 & 4 & 8 & 3 & 6 & 12 & 11 & 9 & 5 & 10 & 7
\end{array}
\]

\[
\begin{array}{c|cccccccccccc}
N & 1 & 2 & 3 & 4 & 5 & 6 & 7 & 8 & 9 & 10 & 11 & 12 \\ \hline
I & 0 & 1 & 4 & 2 & 9 & 5 & 11 & 3 & 8 & 10 & 7 & 6
\end{array}.
\]

Im Canon Arithmeticus von \gesperrt{Jacobi} (1839) sind für alle
Primzahlen im ersten Tausend Indextabellen zusammengestellt. In
kleinerem Umfange nach Berechnungen von \gesperrt{Ostrogradsky} in der „Theorie
der Congruenzen“ von \gesperrt{Tschebyscheff} (deutsch von
\gesperrt{Schapira}, Berlin 1889).

Es ist noch zu bemerken, dass der Index von 1 immer gleich 0 ist, und
dass der Index von $-1$ (oder von $n-1$) immer gleich $\dfrac{n-1}{2}$
ist. Denn da $(-1)^2$ den Index 0 oder $n-1$ hat, so muss $-1$, da es
nicht den Index 0 hat, und das Doppelte seines Index gleich 0 oder
$n-1$ ist, den Index $\dfrac{n-1}{2}$ haben; in Formeln:
\weq{21}{\operatorname{ind} 1 \equiv 0, \quad \operatorname{ind}(-1) \equiv \frac{n-1}{2} \pmod{n-1}.}

\begin{center}
\S.\,137.

\gesperrt{Multiplication und Theilung der trigonometrischen Functionen.}
\end{center}

Wir haben schon auf den Zusammenhang hingewiesen, der zwischen der
Theorie der Einheitswurzeln und der geometrischen Aufgabe der Theilung
der Kreisperipherie besteht. Die Verallgemeinerung dieser Aufgabe
führt auf die Theilung eines beliebigen Winkels. Auch diese Aufgabe
führt auf algebraische Gleichungen, von denen wir die für die
Dreitheilung schon früher für die allgemeine Auflösung der cubischen
Gleichungen benutzt haben (\S.\,112).

Wir beginnen jetzt mit der Aufstellung der allgemeinen Formeln.

Nach dem \gesperrt{Moivre}'schen Satze (\S.\,33) ist, wenn $n$ eine
beliebige positive ganze Zahl und $\varphi$ ein beliebiger Winkel ist,
%CONT>>
